
% =========================================================
\section{Problem setting and standing hypotheses}
\label{sec:setting}
% =========================================================

Throughout the paper, $\langle\cdot,\cdot\rangle$ and
$\|\cdot\|_2$ denote the standard Euclidean inner product and norm
on $\mathbb{R}^d$, respectively. All gradients, Hessians, and spectral
quantities are understood with respect to this Euclidean structure.

Let $\mathcal U\subset\mathbb{R}^d$ be an open set and let
\[
V:\mathcal U\to\mathbb{R}
\]
be a smooth potential. For a step size $\eta>0$, the associated
gradient-descent map is
\begin{equation}
G_\eta(x)
=
x-\eta\nabla V(x).
\label{eq:smooth-gd-map}
\end{equation}

Our aim is to study how a smooth flip bifurcation of
\eqref{eq:smooth-gd-map} interacts with the switching boundary
generated intrinsically by gradient clipping.

% ---------------------------------------------------------
\subsection{Clipped gradient maps}
\label{subsec:clipped-map}
% ---------------------------------------------------------

Let $\|\cdot\|_\diamond$ be a norm on $\mathbb{R}^d$.
For a clipping threshold $c>0$, define
\begin{equation}
\mathcal C_c^\diamond(g)
=
\begin{cases}
g,
&
\|g\|_\diamond\leq c,
\\[1mm]
\dfrac{c}{\|g\|_\diamond}\,g,
&
\|g\|_\diamond>c,
\end{cases}
\qquad
g\in\mathbb{R}^d.
\label{eq:clipping-operator}
\end{equation}
Equivalently,
\[
\mathcal C_c^\diamond(g)
=
\min\left\{
1,\frac{c}{\|g\|_\diamond}
\right\}g,
\]
with the convention $\mathcal C_c^\diamond(0)=0$.

The clipped gradient map is
\begin{equation}
F_{\eta,c}(x)
=
x-\eta\mathcal C_c^\diamond(\nabla V(x)).
\label{eq:clipped-gradient-map}
\end{equation}

The standard Euclidean gradient-clipping rule corresponds to
\[
\|\cdot\|_\diamond=\|\cdot\|_2.
\]

Associated with \eqref{eq:clipped-gradient-map} is the endogenous
switching hypersurface
\begin{equation}
\Sigma_c
=
\left\{
x\in\mathcal U:
\|\nabla V(x)\|_\diamond=c
\right\}.
\label{eq:switching-hypersurface}
\end{equation}
It separates the local phase space into the smooth region
\begin{equation}
\Omega_c^-
=
\left\{
x\in\mathcal U:
\|\nabla V(x)\|_\diamond<c
\right\}
\label{eq:smooth-region}
\end{equation}
and the clipped region
\begin{equation}
\Omega_c^+
=
\left\{
x\in\mathcal U:
\|\nabla V(x)\|_\diamond>c
\right\}.
\label{eq:clipped-region}
\end{equation}

Accordingly,
\begin{equation}
F_{\eta,c}(x)
=
G_\eta(x)
=
x-\eta\nabla V(x),
\qquad
x\in\Omega_c^-,
\label{eq:smooth-branch}
\end{equation}
whereas in $\Omega_c^+$,
\begin{equation}
F_{\eta,c}(x)
=
x
-
\eta c
\frac{\nabla V(x)}
{\|\nabla V(x)\|_\diamond}.
\label{eq:clipped-branch}
\end{equation}

Thus $F_{\eta,c}$ is continuous and piecewise smooth, while its
switching geometry is not prescribed independently of the dynamics:
the boundary $\Sigma_c$ is generated by the same potential $V$
that defines the gradient map.

% ---------------------------------------------------------
\subsection{The distinguished critical point}
\label{subsec:critical-point}
% ---------------------------------------------------------

We focus on the dynamics near an isolated nondegenerate local
minimum $x_*$ of $V$.

Set
\begin{equation}
H_*:=D^2V(x_*).
\label{eq:H-star}
\end{equation}
Since $H_*$ is symmetric, choose an orthonormal eigenbasis
$\{v_i\}_{i=1}^d$ satisfying
\begin{equation}
H_*v_i=\lambda_i v_i,
\qquad
0<\lambda_1\leq\lambda_2\leq\cdots\leq\lambda_d.
\label{eq:H-eigendecomposition}
\end{equation}

We denote the maximal eigenvalue and its corresponding unit
eigenvector by
\begin{equation}
\lambda_*:=\lambda_d,
\qquad
v_*:=v_d.
\label{eq:critical-eigenpair}
\end{equation}

The linearization of the smooth gradient map at $x_*$ is
\begin{equation}
DG_\eta(x_*)
=
I-\eta H_*,
\label{eq:GD-linearization}
\end{equation}
and hence its multipliers are
\begin{equation}
\rho_i(\eta)
=
1-\eta\lambda_i,
\qquad
i=1,\ldots,d.
\label{eq:GD-multipliers}
\end{equation}

The distinguished step size at which the maximal-curvature
multiplier reaches $-1$ is
\begin{equation}
\eta_f
:=
\frac{2}{\lambda_*}.
\label{eq:eta-f}
\end{equation}

At $\eta=\eta_f$,
\begin{equation}
\rho_d(\eta_f)=-1.
\label{eq:critical-multiplier}
\end{equation}

The parameter
\begin{equation}
\mu:=\eta-\eta_f
\label{eq:mu-definition}
\end{equation}
will be used as the local bifurcation parameter throughout the paper.

% ---------------------------------------------------------
\subsection{Higher-order tensors and the flip coefficient}
\label{subsec:flip-coefficient-notation}
% ---------------------------------------------------------

We write
\begin{equation}
T_*:=D^3V(x_*),
\qquad
Q_*:=D^4V(x_*),
\label{eq:TQ-definitions}
\end{equation}
viewed as symmetric multilinear forms.

For $u,w\in\mathbb{R}^d$, let
\begin{equation}
T_*[u,w,\cdot]^\sharp
\in\mathbb{R}^d
\label{eq:sharp-definition}
\end{equation}
denote the Euclidean Riesz representative of the linear functional
$z\mapsto T_*[u,w,z]$; namely,
\begin{equation}
\left\langle
T_*[u,w,\cdot]^\sharp,z
\right\rangle
=
T_*[u,w,z]
\qquad
\text{for all }z\in\mathbb{R}^d.
\label{eq:sharp-identity}
\end{equation}

Define
\begin{equation}
g_*
:=
T_*[v_*,v_*,\cdot]^\sharp
\label{eq:g-star}
\end{equation}
and
\begin{equation}
h_*
:=
H_*^{-1}g_*.
\label{eq:h-star}
\end{equation}

The nonlinear quantity governing the local flip bifurcation is
\begin{equation}
\Gamma_*
:=
3T_*[v_*,v_*,h_*]
-
Q_*[v_*,v_*,v_*,v_*].
\label{eq:Gamma-star}
\end{equation}

Equivalently, using the eigenbasis
\eqref{eq:H-eigendecomposition},
\begin{equation}
\Gamma_*
=
3
\sum_{i=1}^d
\frac{
T_*[v_*,v_*,v_i]^2
}{
\lambda_i
}
-
Q_*[v_*,v_*,v_*,v_*].
\label{eq:Gamma-eigenbasis}
\end{equation}

With the flip-normal-form convention used in this paper, define
\begin{equation}
\ell_*
:=
\frac{\eta_f}{6}\Gamma_*
=
\frac{\Gamma_*}{3\lambda_*}.
\label{eq:ell-star}
\end{equation}
Thus $\ell_*>0$ corresponds to the supercritical case in which a
stable small-amplitude period-two branch is created for
$\eta>\eta_f$.

The expression \eqref{eq:Gamma-eigenbasis} is the multivariate
gradient-map form of the nonlinear flip criterion: the cubic
coupling of the critical direction to all Hessian eigendirections
contributes to the sign of the bifurcation coefficient, rather than
only the derivatives taken purely along $v_*$.

% ---------------------------------------------------------
\subsection{Standing hypotheses}
\label{subsec:standing-hypotheses}
% ---------------------------------------------------------

The following assumptions are imposed unless stated otherwise.

\begin{enumerate}[
label=\textbf{(H\arabic*)},
leftmargin=2.8em,
itemsep=0.8em
]

\item
\label{H1}
\textbf{Regularity.}
The potential satisfies
\[
V\in C^7(\mathcal U,\mathbb{R}).
\]

\item
\label{H2}
\textbf{Nondegenerate strict local minimum.}
The distinguished point $x_*\in\mathcal U$ satisfies
\begin{equation}
\nabla V(x_*)=0,
\qquad
H_*=D^2V(x_*)\succ0.
\label{eq:H2}
\end{equation}

\item
\label{H3}
\textbf{Simple maximal curvature.}
The largest Hessian eigenvalue is simple:
\begin{equation}
\lambda_{d-1}<\lambda_d=\lambda_*.
\label{eq:H3}
\end{equation}

\item
\label{H4}
\textbf{Supercritical flip nondegeneracy.}
The coefficient defined in \eqref{eq:Gamma-star} satisfies
\begin{equation}
\Gamma_*>0,
\label{eq:H4}
\end{equation}
or equivalently
\[
\ell_*>0.
\]

\item
\label{H5}
\textbf{Smooth clipping geometry.}
The clipping norm $\|\cdot\|_\diamond$ is of class $C^3$ on
$\mathbb{R}^d\setminus\{0\}$.
\end{enumerate}

% ---------------------------------------------------------
\subsection{Immediate spectral consequences}
\label{subsec:spectral-consequences}
% ---------------------------------------------------------

Assumptions \textbf{(H2)--(H3)} imply that the spectral hypotheses
required for a simple codimension-one flip bifurcation are automatic
for the gradient map.

Indeed, for $i=1,\ldots,d-1$,
\[
0<\lambda_i<\lambda_*,
\]
and therefore, at $\eta=\eta_f$,
\begin{equation}
-1
<
1-\frac{2\lambda_i}{\lambda_*}
<
1.
\label{eq:noncritical-spectrum}
\end{equation}
Hence $-1$ is the unique multiplier of $DG_{\eta_f}(x_*)$ on the
unit circle.

Moreover,
\begin{equation}
\frac{d}{d\eta}\rho_d(\eta)
=
-\lambda_*,
\label{eq:eigenvalue-crossing}
\end{equation}
so that
\begin{equation}
\left.
\frac{d}{d\eta}\rho_d(\eta)
\right|_{\eta=\eta_f}
=
-\lambda_*
\neq0.
\label{eq:transversal-crossing}
\end{equation}
Thus the critical multiplier crosses $-1$ transversally as $\eta$
passes through $\eta_f$.

Thus, under the spectral assumptions \textbf{(H2)--(H3)}, the sign of
$\Gamma_*$ provides the remaining nonlinear nondegeneracy condition
that distinguishes the supercritical and subcritical flip cases.

% ---------------------------------------------------------
\subsection{Local form of the switching geometry}
\label{subsec:local-switching-preview}
% ---------------------------------------------------------

We record one elementary observation that motivates the later
switching analysis.

Let
\[
x=x_*+sv_*+r(s),
\qquad
r(s)=O(s^2),
\qquad
s\to0.
\]
Since $\nabla V(x_*)=0$ and $H_*v_*=\lambda_*v_*$,
Taylor expansion gives
\begin{equation}
\nabla V(x)
=
\lambda_*s\,v_*
+
O(s^2).
\label{eq:gradient-critical-expansion}
\end{equation}

By positive homogeneity and smoothness of the clipping norm away
from the origin,
\begin{equation}
\|\nabla V(x)\|_\diamond
=
\lambda_*
\|v_*\|_\diamond
|s|
+
O(s^2).
\label{eq:norm-critical-expansion}
\end{equation}

In particular, the leading-order displacement of the switching
boundary from $x_*$ in the critical direction is nondegenerate.

For the Euclidean clipping rule,
$\|v_*\|_2=1$, and
\eqref{eq:norm-critical-expansion} reduces to
\begin{equation}
\|\nabla V(x)\|_2
=
\lambda_*|s|
+
O(s^2).
\label{eq:euclidean-critical-expansion}
\end{equation}

 A precise expansion of $\Sigma_c$ along the center manifold will be
derived later in Section~\ref{sec:switching-geometry}.

% ---------------------------------------------------------
\subsection{Scope of the present paper}
\label{subsec:scope}
% ---------------------------------------------------------

The positive-definiteness assumption in \textbf{(H2)} deliberately
restricts the present work to an isolated strict local minimum.
This excludes the additional unit multiplier $+1$ that occurs when
the Hessian has a nontrivial kernel, as in Morse--Bott manifolds of
minimizers.

The purpose of this restriction is to isolate the interaction between
a simple smooth flip bifurcation and the endogenous switching
boundary generated by clipping.  Extensions involving manifolds of
minimizers, multiple critical multipliers, or degenerate flip
coefficients require a higher-dimensional center dynamics and are
outside the scope of the present paper.

